\documentclass[11pt]{article}
\usepackage[margin=1in]{geometry}
\usepackage{amsmath,amssymb,amsthm}
\usepackage{mathtools}
\usepackage{enumitem}
\usepackage{hyperref}
\usepackage{cleveref}
\usepackage{booktabs}
\usepackage{tabularx}
\usepackage{microtype}
\usepackage{cite}
\usepackage{xurl}
\usepackage{path}
\hypersetup{colorlinks=true,linkcolor=blue,citecolor=blue,urlcolor=blue}
\newtheorem{definition}{Definition}
\newtheorem{lemma}{Lemma}
\title{Challenge Answer Review Menus for Successor Maintenance\\\large Smallest-Sufficient Referee Handoffs from Dockets, Profiles, Packets, and Ladders}
\author{Anonymous}
\date{Unified archive note v0.15 (2026-03-28; archive v727)}
\begin{document}
\maketitle

\begin{abstract}
The late-stage answer side is now rich enough that a brief paper faces one last repetition problem: when a referee asks a concrete follow-up question, which maintained object is the smallest sufficient thing to hand over? The answer is not always ``the full packet,'' and it is not always ``just cite the docket.'' This note isolates that response-selection problem. It defines a compact \emph{challenge answer review menu}: one audience/question map from common referee request classes to the smallest sufficient already-maintained answer companion. The menu imports publication profiles for inline checks, citation dockets for exact-location and full-audit checks, disclosure packets for requestable archive handoffs, and escalation ladders for reveal-order checks. It does not create a new semantic answer layer. It just keeps referee follow-up policy stable and auditable under space pressure.
\end{abstract}

\paragraph{Non-redundancy note.}
Synthesis~50 owns smallest-sufficient answer citation dockets, Synthesis~51 owns answer publication profiles, Synthesis~52 owns answer disclosure packets, Synthesis~53 owns answer escalation ladders, Synthesis~31 owns the downstream suffix, Synthesis~32 owns the stable release spine, Synthesis~55 owns the theorem-to-review bridge across those layers, Synthesis~17 owns the concrete worked instantiation, and Synthesis~76 owns the wider theorem-to-publication route that stops before this answer-side terminal owner.\cite{series:challengeanswercitationdockets,series:challengeanswerpublicationprofiles,series:challengeanswerdisclosurepackets,series:challengeanswerescalationladders,series:downstreamnormalforms,series:releasespine,series:seriesspines,series:workedexample,series:theoremtopublicationspine}
This note owns only the request-class selection among those already-maintained companions. The repeated five reveal-step labels inside the imported escalation ladder remain ladder-owned vocabulary rather than a new review-menu-owned step policy. If a later paper needs the exact paper-facing packet/menu contract behind a source-only public sentence, that aggregation belongs in Synthesis~56 rather than here. If the later paper also needs a stable name for the omitted-support family behind that request-packet handoff, that vocabulary belongs in Synthesis~57 rather than in the menu. If it needs the exact completion criterion for when that promised cut has been fully satisfied, that belongs in Synthesis~58 rather than in the menu.\cite{series:publicrequestcontracts,series:requestableevidenceclasses,series:requestfulfillmentcertificates}

\paragraph{Scope.}
This is not a new answer card, carrier slot, audit trail, citation docket, publication profile, disclosure packet, escalation ladder, or theorem-to-publication route.
It is a thin review-time policy note saying which imported object should be handed over for one concrete referee-style follow-up class.

\paragraph{Later-terse-paper citation posture.}
Later terse papers should cite Synthesis~54 only when the live issue is the answer-side request-class selection itself --- which maintained companion answers \textsf{brief-inline}, \textsf{expanded-inline}, \textsf{exact-location}, \textsf{full-audit}, \textsf{request-packet}, or \textsf{reveal-order} for one audience/question pair --- or when they need the compact answer-side terminal citation normal form or its one-clause sentence rendering exported below. If the live issue is instead the exact docket entry, publication-profile floor, disclosure-packet contents, or escalation-order policy, stop at Synthesis~50--53. If the question has already reduced to downstream summary / support / citation / quote / clause reuse, stop at Synthesis~31; if it has widened to release-stage routing, stop at Synthesis~32. Once a checked-answer handoff already exists and the only remaining issue is whether a later terse paper may stop there, must widen to the adjacent request-side capstone, or must reopen a named owner, stop at Synthesis~73 for the maintained request-side terminal stop, drop to the worked publication-boundary card in Synthesis~17 when one concrete checked-answer/request-tail seam is wanted, and widen to Synthesis~74--75 only when the live issue is the abstract paired-tail discipline or abstract stop/widen/reopen judgment rather than that concrete seam.\cite{series:challengeanswercitationdockets,series:challengeanswerpublicationprofiles,series:challengeanswerdisclosurepackets,series:challengeanswerescalationladders,series:downstreamnormalforms,series:releasespine,series:workedexample,series:publicrequestresponsepacketrefreshresponsepacketclosureverdicts,series:pairedterminalcitationdiscipline,series:pairedterminalcutoverrules}

\paragraph{One branched answer-terminal route.}
Later terse papers should usually treat the late answer side as one branched terminal route rather than as a fake linear stack: Synthesis~51 owns the inline-budget branch, Synthesis~50 owns the exact-location/full-audit citation branch, Synthesis~52 owns the request-packet branch, Synthesis~53 owns the reveal-order branch, and Synthesis~54 is the capstone request-class selector that chooses among those already-owned branches for one audience/question pair. This note is therefore the capstone rather than another branch. Stop at Synthesis~50--53 whenever one branch owner is still genuinely live; stop here only when the branch owners are already fixed and the remaining issue is which concrete referee request class should receive which maintained companion; then stop at Synthesis~73 for the maintained request-side terminal stop, drop to the worked publication-boundary card in Synthesis~17 when one concrete checked-answer/request-tail seam is wanted rather than another owner summary, and widen to Synthesis~74--75 only when the remaining issue is the abstract paired-tail discipline or abstract fail-closed cutover judgment; fail closed back to Synthesis~76 only when the real issue is no longer this terminal owner but which wider theorem/worked/release/notary/claim/release owner should carry the claim next.

\paragraph{Exports.}
This note exports (i) challenge answer review menus, (ii) a request-class coverage rule, (iii) a smallest-sufficient review-response rule, (iv) an imported-locality rule saying that later papers may cite one maintained menu when they need a stable answer to ``what exactly should I hand over for this kind of follow-up request?'', (v) a shared-review-vocabulary rule saying that repeated request-class labels may travel through later routes or spine mirrors only as imported menu vocabulary, (vi) an imported-escalation-vocabulary rule saying that repeated reveal-step labels mirrored from the imported escalation ladder remain ladder-owned vocabulary rather than a new menu-owned reveal policy, and (vii) an answer-side terminal sentence normal form that renders the already-owned four-anchor checked-answer stop as one exact one-clause sentence for later terse papers.

\section{Why answer review menus deserve their own note}
Escalation ladders already say in what order support forms should be revealed.
Disclosure packets already say what exact archive slice should be handed over once the reader has outgrown inline, exact-location, and full-audit support.
But later papers still face a narrower operational question: if a referee asks for one concrete kind of follow-up, what is the smallest sufficient maintained object to hand over right now?
Without a maintained answer-review menu, terse papers either overshare the whole packet or underspecify the response in prose.
So the archive benefits from one small note that turns common request classes into explicit handoff choices.

\begin{definition}[Challenge answer review menu]
Fix one concrete base$\to$successor comparison, one audience key $a$, and one downstream question key $q$.
A \emph{challenge answer review menu} is a tuple
\[
M(a,q)=(a,q,P,C,D,L,R)
\]
containing one imported answer-publication-profile key $P$, one imported answer-citation-docket key $C$, one imported answer-disclosure-packet key $D$, one imported answer-escalation-ladder key $L$, and an ordered family $R=(r_1,\dots,r_6)$ of review-response entries for the request classes
\[
\textsf{brief-inline},\ \textsf{expanded-inline},\ \textsf{exact-location},\ \textsf{full-audit},\ \textsf{request-packet},\ \textsf{reveal-order}.
\]
Each entry names the smallest sufficient already-maintained object for that request class together with its canonical citation path and the reason that object is sufficient.
\end{definition}

\begin{definition}[Request-class coverage rule]
A challenge answer review menu satisfies the \emph{request-class coverage rule} if it covers exactly the six request classes above, in that fixed order, for the same audience/question pair.
So the menu cannot silently omit the common referee follow-ups that brief papers repeatedly encounter.
\end{definition}

\begin{definition}[Smallest-sufficient review-response rule]
A challenge answer review menu satisfies the \emph{smallest-sufficient review-response rule} if it chooses:
\begin{enumerate}[leftmargin=*]
\item the imported publication profile for both brief-inline and expanded-inline checks,
\item the imported citation docket for exact-location and full-audit checks,
\item the imported disclosure packet for request-packet handoff, and
\item the imported escalation ladder for reveal-order checks.
\end{enumerate}
So the menu changes only request-class selection; it does not widen any imported answer companion.
\end{definition}

\begin{definition}[Imported-locality rule]
A challenge answer review menu satisfies the \emph{imported-locality rule} if every imported key belongs to the same audience/question pair and every review entry copies the imported inline form, docket entry, disclosure-packet contents, or reveal-order steps exactly.
So the menu may not create new answer text, new carrier locations, or new audit paths while pretending only to select a handoff object.
\end{definition}

\begin{definition}[Shared-review-vocabulary rule]
A challenge answer review menu satisfies the \emph{shared-review-vocabulary rule} if its repeated request-class labels may recur in later routes or series-spine mirrors only as imported review-menu vocabulary: those later objects may mirror that \textsf{brief-inline}/\textsf{expanded-inline}/\textsf{exact-location}/\textsf{full-audit}/\textsf{request-packet}/\textsf{reveal-order} ladder, but they do not thereby become the owner of the concrete handoff choice for any label.
So later terse notes may reuse the six class names without silently re-owning the audience/question-scoped menu selection that picked the maintained object for each one.
\end{definition}

\begin{definition}[Imported-escalation-vocabulary rule]
A challenge answer review menu satisfies the \emph{imported-escalation-vocabulary rule} if any repeated reveal-step labels that appear through its imported escalation-ladder entry remain imported ladder vocabulary: the review menu may point at \textsf{default-inline}, \textsf{expanded-inline}, \textsf{exact-location}, \textsf{full-audit}, and \textsf{request-packet} as the maintained reveal order, but it does not thereby become the owner of that step order.
So the review menu chooses which companion to hand over for \textsf{reveal-order} checks, while the ladder still owns the internal five-step reveal policy it imports.
\end{definition}


\begin{definition}[Answer-side terminal citation normal form]
Fix one audience/question answer whose review menu satisfies the request-class coverage rule, the smallest-sufficient review-response rule, the imported-locality rule, the shared-review-vocabulary rule, and the imported-escalation-vocabulary rule.
Its \emph{answer-side terminal citation normal form} is the four-anchor route
\[
\textsf{public-anchor} \to \textsf{checked-answer-card} \to \textsf{answer-disclosure-packet} \to \textsf{answer-review-menu},
\]
where the checked-answer card carries the default inline publication profile for that audience/question answer and the disclosure packet and review menu import the exact-location, full-audit, and reveal-order companions that may be reopened later.
So the normal form records the compared release pair, the exact checked sentence, the exact answer-side omitted-support cut, and the terminal referee-handoff owner without rewalking every intermediate answer companion each time.
\end{definition}

\begin{definition}[Answer-side terminal sentence normal form]
An answer-side terminal citation normal form exports an \emph{answer-side terminal sentence normal form} if it supplies exactly one clause:
\begin{quote}\small
For the checked-answer path, stop at: \textsf{public-anchor} $\to$ \textsf{checked-answer-card} $\to$ \textsf{answer-disclosure-packet} $\to$ \textsf{answer-review-menu}.
\end{quote}
So the sentence form adds no new owner and no new semantic field: it is only the stable one-clause rendering of the already-owned four-anchor answer-side stop.
\end{definition}

\begin{lemma}[Later terse papers may cite the answer-side terminal citation normal form]
Assume one audience/question answer has a challenge answer review menu satisfying the request-class coverage rule, the smallest-sufficient review-response rule, the imported-locality rule, the shared-review-vocabulary rule, and the imported-escalation-vocabulary rule.
Then a later terse paper may cite that answer's answer-side terminal citation normal form instead of reopening the whole checked-answer chain, provided the later paper is not trying to re-litigate exact-location, full-audit, or reveal-order content.
\end{lemma}

\begin{proof}
The checked-answer-card anchor fixes the exact emitted sentence together with its default inline publication profile.
The disclosure-packet anchor fixes the exact answer-side omitted-support cut and inherits the exact-location and full-audit companions it names through the imported citation docket.
The review-menu anchor fixes the smallest-sufficient referee handoff object for each request class and imports the escalation ladder when reveal order itself becomes the question.
So the four-anchor route preserves the beginning, terminal answer-side cut, and terminal handoff owner of the checked-answer path while leaving the omitted publication-profile, citation-docket, carrier-slot, audit-trail, and escalation-ladder details auditable on demand rather than mandatory in every later citation. Once that four-anchor route is fixed, the answer-side terminal sentence normal form is simply its stable one-clause rendering for later terse papers rather than a new owner layer.\qedhere
\end{proof}

\begin{lemma}[Review menus keep referee handoffs brief without changing answer ownership]
Assume every challenge answer review menu satisfies the request-class coverage rule, the smallest-sufficient review-response rule, the imported-locality rule, the shared-review-vocabulary rule, and the imported-escalation-vocabulary rule.
Then a later paper may cite one menu as the maintained guide for referee follow-up handling on one audience/question answer without changing the ownership of inline publication choices, exact-location support, full-audit support, requestable archive slices, or reveal order.
\end{lemma}

\begin{proof}
The request-class coverage rule ensures that the common referee requests are covered explicitly.
The smallest-sufficient review-response rule ensures that each request class resolves to one already-owned object and no larger one.
The imported-locality rule then prevents the menu from rewriting any imported answer companion. The shared-review-vocabulary rule adds the later-seam guarantee: even when the same six request-class names are mirrored into route catalogs or series-spine surfaces, those mirrors still import vocabulary from the menu rather than re-owning the actual handoff choice for a given audience/question pair. The imported-escalation-vocabulary rule adds the upstream seam guarantee: even when later surfaces repeat the five reveal-step labels through the imported escalation ladder, those labels still travel as ladder-owned vocabulary rather than as a new menu-owned reveal policy.
So the menu changes only policy packaging: it stabilizes what to hand over for one review request class, but it does not alter the meaning or ownership of any imported rung.\qedhere
\end{proof}

\begin{table}[t]
\centering
\small
\begin{tabularx}{\linewidth}{@{}l l X@{}}
\toprule
Request class & Canonical object & Referee question answered \\
\midrule
Brief inline check & publication profile & What is the smallest inline answer form the terse paper should print by default? \\
Expanded inline check & publication profile & What larger inline provenance form stays on the same emitted answer? \\
Exact-location check & citation docket & Which exact shipped field or reusable clause carries that answer? \\
Full-audit check & citation docket & Which full leftward walk discharges the answer to narrow terminal owners? \\
Request-packet check & disclosure packet & Which exact archive slice and validator companions should be handed over on request? \\
Reveal-order check & escalation ladder & In what maintained order should those support forms be revealed? \\
\bottomrule
\end{tabularx}
\caption{A review menu does not add another answer stage. It chooses the smallest sufficient already-owned companion for one referee follow-up class.}
\label{tab:answerreviewmenu}
\end{table}

\paragraph{A minimal public answer-review-menu card.}
\begin{table}[t]
\centering
\small
\begin{tabularx}{\linewidth}{@{}lX@{}}
\toprule
Field & Worked answer \\
\midrule
Release-note status menu & \nolinkurl{release_note_public_status_sentence_answer_review_menu} binds audience \texttt{release\_note} and question \texttt{public\_status\_sentence} to the imported profile \nolinkurl{release_note_public_status_sentence_answer_publication_profile}, docket \nolinkurl{release_note_public_status_sentence_answer_citation_docket}, disclosure packet \nolinkurl{release_note_public_status_sentence_answer_disclosure_packet}, and escalation ladder \nolinkurl{release_note_public_status_sentence_answer_escalation_ladder} \\
Auditor compact-diff menu & \nolinkurl{auditor_compact_diff_summary_answer_review_menu} binds audience \texttt{auditor} and question \texttt{compact\_diff\_summary} to the imported profile \nolinkurl{auditor_compact_diff_summary_answer_publication_profile}, docket \nolinkurl{auditor_compact_diff_summary_answer_citation_docket}, disclosure packet \nolinkurl{auditor_compact_diff_summary_answer_disclosure_packet}, and escalation ladder \nolinkurl{auditor_compact_diff_summary_answer_escalation_ladder} \\
Inline floor & Both menus resolve \nolinkurl{brief_inline_check} and \nolinkurl{expanded_inline_check} through their imported publication profiles to the same checked answer card, changing only the declared inline-provenance mode rather than the emitted answer \\
Audit floor & Both menus resolve \nolinkurl{exact_location_check} and \nolinkurl{full_audit_check} through their imported citation dockets to the exact carrier slot and leftward audit trail already maintained for that same answer \\
Packet / ladder floor & Both menus resolve \nolinkurl{request_packet_check} to the imported disclosure packet and \nolinkurl{reveal_order_check} to the imported escalation ladder, so the menu chooses the handoff object without re-owning packet contents or reveal-order semantics \\
Escalation pointer & Stop at Synthesis~50--53 when the live issue is the narrower profile / docket / packet / ladder owner; stop here when the live issue is which maintained companion one request class selects; stop at Synthesis~73 for the maintained request-side terminal stop when that is the live issue; drop to the worked publication-boundary card in Synthesis~17 when one concrete checked-answer/request-tail seam is wanted; widen to Synthesis~74--75 only when the remaining issue is the abstract paired-tail discipline or abstract fail-closed cutover rather than the answer-side handoff itself \\
\bottomrule
\end{tabularx}
\caption{A minimal public answer-review-menu card. This is the smallest answer-side public answer later terse papers can quote directly when the live issue is which already-maintained companion should answer one concrete referee follow-up, rather than the narrower imported owner or the wider paired-terminal cutover.}
\label{tab:minimal-public-answer-review-menu-card}
\end{table}

This card is intentionally non-decorative: it pins the actual worked menu keys, the imported companion owners they bind, and the answer-side stop rule without silently re-owning inline publication choices, exact-location entries, full-audit walks, requestable archive slices, or reveal-order policy.
\cite{series:workedexample}

\paragraph{A tiny answer-side terminal sentence card.}
\begin{table}[t]
\centering
\small
\begin{tabularx}{\linewidth}{@{}lX@{}}
\toprule
Field & Minimal public answer \\
\midrule
Checked-answer clause floor & \texttt{For the checked-answer path, stop at:}\newline
\texttt{\nolinkurl{example_compare_report.json}}\newline
$\to$\newline
\texttt{\nolinkurl{auditor_compact_diff_summary_answer_card}}\newline
$\to$\newline
\texttt{\nolinkurl{auditor_compact_diff_summary_answer_disclosure_packet}}\newline
$\to$\newline
\texttt{\nolinkurl{auditor_compact_diff_summary_answer_review_menu}}. \\
Sentence-owner floor & This one clause is only the sentence rendering of the already-owned answer-side terminal citation normal form; the compared release pair, exact checked sentence, omitted-support cut, and terminal referee handoff owner are still carried by the same four anchors \\
Imported-companion floor & The publication profile, citation docket, carrier slot, audit trail, and escalation ladder remain auditable support behind the clause and reopen only when exact-location, full-audit, or reveal-order content is itself the live question \\
Staleness trigger floor & This clause goes stale as soon as the checked answer card, the disclosure-packet cut, the review-menu key, or the compared public anchor changes, even if a nearby later paper is discussing a similar answer \\
Escalation pointer & Stop at Synthesis~50--53 when the narrower imported owner is still live; stop here when one exact answer-side checked-answer sentence stop is the live issue; stop at Synthesis~73 for the maintained request-side terminal stop when that is the live issue; drop to the worked publication-boundary card in Synthesis~17 when one concrete checked-answer/request-tail seam is wanted; widen to Synthesis~74--75 only when the remaining issue is the abstract question whether a later terse paper may stay on that answer-side stop or must cross the adjacent request-side cut \\
\bottomrule
\end{tabularx}
\caption{A tiny answer-side terminal sentence card. This is the smallest public sentence-level rendering later terse papers can quote directly when the live issue is the exact checked-answer stop itself rather than the narrower imported owners or the wider paired-terminal cutover.}
\label{tab:tiny-answer-side-terminal-sentence-card}
\end{table}

This sentence card is intentionally non-decorative: it gives later terse papers one exact checked-answer stop clause without silently re-owning the imported profile, docket, packet, or ladder layers that still supply the auditable middle support.\cite{series:workedexample}

\paragraph{A tiny checked-answer stop fill.}
For the worked auditor compact-diff answer, later terse papers may now write one exact clause:
\begin{quote}\small
For the checked-answer path, stop at: \nolinkurl{example_compare_report.json} $\to$ \nolinkurl{auditor_compact_diff_summary_answer_card} $\to$ \nolinkurl{auditor_compact_diff_summary_answer_disclosure_packet} $\to$ \nolinkurl{auditor_compact_diff_summary_answer_review_menu}.
\end{quote}
That clause stays honest only while the same compared public anchor, checked answer card, disclosure-packet cut, and review-menu key remain live. If the live issue is instead the exact carrier slot, the full leftward audit, or the reveal-order rung, the paper should reopen Synthesis~50--53 rather than keep quoting the clause; if the answer-side stop is already fixed and the real question is the maintained request-side terminal stop, it should stop first at Synthesis~73; if the real question is one concrete checked-answer/request-tail seam, it should drop first to the worked publication-boundary card in Synthesis~17; only if the live issue is the abstract paired-tail discipline or abstract stop/widen/reopen judgment should it leave this note for Synthesis~74--75.\cite{series:workedexample}

\paragraph{A forced public answer-review-menu matrix.}
\begin{table}[t]
\centering
\small
\begin{tabularx}{\linewidth}{@{}lX@{}}
\toprule
Field & Forced public answer \\
\midrule
Release-note menu-entry floor & The pair \nolinkurl{release_note_public_status_sentence_answer_review_menu} plus \nolinkurl{exact_location_check} still forces the single selected owner \nolinkurl{release_note_public_status_sentence_answer_carrier_slot}; the same menu plus \nolinkurl{full_audit_check} still forces \nolinkurl{release_note_public_status_sentence_answer_audit_trail}; and the same menu plus \nolinkurl{request_packet_check} still forces the disclosure packet \nolinkurl{release_note_public_status_sentence_answer_disclosure_packet} \\
Auditor menu-entry floor & The pair \nolinkurl{auditor_compact_diff_summary_answer_review_menu} plus \nolinkurl{brief_inline_check} or \nolinkurl{expanded_inline_check} still forces the checked answer card \nolinkurl{auditor_compact_diff_summary_answer_card} under the corresponding inline mode; the same menu plus \nolinkurl{exact_location_check} still forces \nolinkurl{auditor_compact_diff_summary_answer_carrier_slot}; the same menu plus \nolinkurl{full_audit_check} still forces \nolinkurl{auditor_compact_diff_summary_answer_audit_trail}; and the same menu plus \nolinkurl{request_packet_check} or \nolinkurl{reveal_order_check} still forces the disclosure packet \nolinkurl{auditor_compact_diff_summary_answer_disclosure_packet} or escalation ladder \nolinkurl{auditor_compact_diff_summary_answer_escalation_ladder} \\
Imported-owner floor & Those menu-entry shortcuts remain honest only while the same audience/question pair, the same emitted answer lineage, the same request-class label, and the same imported profile / docket / packet / ladder keys remain bound by the menu; the menu does not get to keep the same answer after any of those imported owners changes \\
Vocabulary floor & Each forced entry still carries the menu-owned review-request vocabulary tag \nolinkurl{imported_vocabulary_menu_owned_selection}, so later terse papers can recover that the class label is shared vocabulary while the concrete handoff is still menu-owned for that audience/question pair \\
Staleness trigger floor & This matrix goes stale as soon as the audience/question pair changes, the emitted answer card changes, a different request class becomes live, the selected object key no longer matches the menu entry, or the menu ceases to expose the same imported companion key beside that entry \\
Escalation pointer & Stop at Synthesis~50--53 when the live issue is still the narrower imported owner; stop here only while one audience/question pair together with one request class still forces the same selected handoff; stop at Synthesis~73 for the maintained request-side terminal stop when that is the live issue; drop to the worked publication-boundary card in Synthesis~17 when one concrete checked-answer/request-tail seam is wanted; widen to Synthesis~74--75 only after the answer-side handoff is fixed and the remaining issue is the abstract question whether a later terse paper may stop there or must cross the adjacent request-side cut \\
\bottomrule
\end{tabularx}
\caption{A forced public answer-review-menu matrix. This is the smallest answer-side public answer later terse papers can quote directly when the live issue is whether one audience/question pair together with one review class still forces the same maintained handoff entry, rather than merely which menu exists in the abstract.}
\label{tab:forced-public-answer-review-menu-matrix}
\end{table}

This matrix is intentionally non-decorative: it pins the actual worked audience/question-class pairs that still force the same answer-side handoff entries, together with the exact conditions under which that forcing remains honest or goes stale, without silently re-owning the imported profile, docket, disclosure-packet, or escalation-ladder objects that neighboring notes already own.\cite{series:workedexample}

\paragraph{One tiny answer-review-menu drill.}
For the worked auditor compact-diff answer, the pair \nolinkurl{auditor_compact_diff_summary_answer_review_menu} plus \nolinkurl{exact_location_check} still forces \nolinkurl{auditor_compact_diff_summary_answer_carrier_slot}. If the live question changes to ``walk every semantic piece leftward,'' the honest stop remains the same menu but the class changes to \nolinkurl{full_audit_check}, so the forced handoff becomes \nolinkurl{auditor_compact_diff_summary_answer_audit_trail}. If the live question instead stays on \nolinkurl{exact_location_check} but the emitted answer card, audience/question pair, or imported citation-docket key changes, the matrix has gone stale and a later terse paper must reopen the exact answer-review-menu owner rather than keep quoting the old carrier-slot entry. Only once the answer-side handoff itself is fixed and the real question becomes the maintained request-side terminal stop should it stop first at Synthesis~73. If the real question becomes one concrete checked-answer/request-tail seam, it should drop first to the worked publication-boundary card in Synthesis~17; only when the live issue is the abstract paired-tail discipline or abstract stop/widen/reopen judgment should it leave this note for Synthesis~74--75.\cite{series:workedexample}

\paragraph{A compact first-reopen-owner card.}
\begin{table}[t]
\centering
\small
\begin{tabularx}{\linewidth}{@{}p{0.31\linewidth}p{0.23\linewidth}X@{}}
\toprule
Failure class & First reopen owner & Why the menu shortcut is no longer honest \\
\midrule
The same audience/question answer still exists but the concrete request-class-to-companion selection changes, one review entry is replaced, or the menu's local handoff sentence is re-cut while the imported companion owners remain nearby & Reopen the answer-review-menu owner in Synthesis~54. & The live issue is still which already-maintained companion should answer one concrete referee follow-up for one audience/question pair, so the capstone request-class selector itself regains control. \\
The brief-inline or expanded-inline entry changes while the rest of the menu remains nearby & Reopen the publication-profile owner in Synthesis~51. & Those two review entries import their actual floor and capsule answers from the publication profile rather than owning inline-budget content locally. Once either inline entry changes, the first honest owner is the imported publication-profile branch. \\
The exact-location or full-audit entry changes while the audience/question answer still looks nearby & Reopen the citation-docket owner in Synthesis~50. & The menu may choose when \textsf{exact-location} or \textsf{full-audit} should be handed over, but it does not own which exact carrier slot or leftward audit trail fills those entries. Once either audit entry changes, the first honest owner is the imported citation-docket branch. \\
The request-packet entry changes while the earlier review classes remain nearby & Reopen the disclosure-packet owner in Synthesis~52. & The menu may route one class to a requestable packet, but it does not own the omitted-support slice itself. Once that handoff object changes, the first honest owner is the imported disclosure-packet branch. \\
The reveal-order entry changes while the answer-side request classes themselves remain nearby & Reopen the escalation-ladder owner in Synthesis~53. & The menu points at one maintained reveal-order companion, but the order and rung content are still owned by the imported escalation ladder. Once the reveal-order entry changes, the first honest owner is the imported reveal-order branch. \\
The audience/question pair, whole-question handoff, or emitted checked-answer lineage changes rather than only one request-class entry & Reopen Synthesis~46 or Synthesis~47 according to whether the first drift is the whole-question bundle or the emitted answer identity. & The menu only stabilizes review-class selection for one already-maintained answer. If the answer-side handoff itself changes, the menu cannot stay honest before the narrower answer-sheet or answer-card owner is revalidated. \\
The answer-side handoff is fixed and the remaining issue is whether a later terse paper may stop there, must cross the adjacent request-side cut, or must reopen a paired-terminal owner & Drop first to \textbf{Synthesis~17} when one concrete checked-answer/request-tail seam is wanted; widen to \textbf{Synthesis~74--75} only for the abstract paired-tail discipline or abstract stop/widen/reopen judgment. & At that point request-class selection is no longer the live question. The paper should leave the answer-side capstone menu, stop first at the worked publication-boundary card when one concrete maintained seam is wanted, and widen to the paired-terminal stop / fail-closed owners only when the live issue is abstract rather than concrete. \\
\bottomrule
\end{tabularx}
\caption{A compact first-reopen-owner card. This is the smallest local map from one stale answer-review-menu shortcut to the first narrower or wider owner that regains authority.}
\label{tab:first-reopen-owner-card-answer-review-menu}
\end{table}

\paragraph{One tiny first-reopen-owner drill.}
If a later terse paper still asks \emph{``which already-maintained companion should answer this same referee follow-up class for this same checked answer?''} but the only changed object is the class-to-companion selection sentence or one review entry itself, Table~\ref{tab:first-reopen-owner-card-answer-review-menu} sends the paper back to Synthesis~54 itself. If the menu is unchanged except that the requested inline entry no longer reuses the same floor or capsule, the same table instead reopens Synthesis~51. If the menu remains fixed and the real question becomes whether the paper may stop at that answer-side handoff or must widen across the adjacent request-side cut, the same table stops first at Synthesis~73 when the maintained request-side terminal stop is the live issue, drops first to the worked publication-boundary card in Synthesis~17 when one concrete checked-answer/request-tail seam is wanted, and widens to Synthesis~74--75 only when the live issue is the abstract paired-tail discipline or abstract stop/widen/reopen judgment rather than that concrete seam.\cite{series:workedexample,series:challengeanswerpublicationprofiles,series:publicrequestresponsepacketrefreshresponsepacketclosureverdicts,series:pairedterminalcitationdiscipline,series:pairedterminalcutoverrules}

\section{Worked example pointer}
The maintained worked bundle now ships \nolinkurl{example_successor_challenge_answer_review_menus.json}.\cite{series:workedexample}
It records four answer review menus for the canned Case-B successor.
Concrete keys already in the ledger include \nolinkurl{release_note_public_status_sentence_answer_review_menu} and \nolinkurl{auditor_compact_diff_summary_answer_review_menu}.
For the release-note public-status menu, the brief and expanded inline request classes resolve to the shared publication profile, the exact-location and full-audit classes resolve to the citation docket entries already selecting the lineage-notice carrier slot and the corresponding audit trail, the request-packet class resolves to the disclosure packet, and the reveal-order class resolves to the escalation ladder for that same answer.
For the auditor compact-diff menu, the same six classes become a concrete stop map rather than a generic taxonomy: \textsf{brief-inline-check} stops at the default-inline answer card, \textsf{expanded-inline-check} at the expanded-inline answer card, \textsf{exact-location-check} at carrier slot \nolinkurl{auditor_compact_diff_summary_answer_carrier_slot}, \textsf{full-audit-check} at audit trail \nolinkurl{auditor_compact_diff_summary_answer_audit_trail}, \textsf{request-packet-check} at disclosure packet \nolinkurl{auditor_compact_diff_summary_answer_disclosure_packet}, and \textsf{reveal-order-check} at escalation ladder \nolinkurl{auditor_compact_diff_summary_answer_escalation_ladder}. Each widens only when the next named need appears: more visible provenance, exact clause location, a contested semantic piece, the omitted archive slice itself, a later paper-facing follow-up class, or a dispute about rung content rather than order.

For later terse papers, those same maintained owners now also support one reusable answer-side terminal citation normal form:
\begin{quote}\small
\nolinkurl{example_compare_report.json}\\
$\to$ \nolinkurl{auditor_compact_diff_summary_answer_card}\\
$\to$ \nolinkurl{auditor_compact_diff_summary_answer_disclosure_packet}\\
$\to$ \nolinkurl{auditor_compact_diff_summary_answer_review_menu}.
\end{quote}
In that shorter citation the checked-answer-card anchor still carries default inline profile \texttt{floor} and visible provenance floor \texttt{base\_release\_id}, \texttt{successor\_release\_id}, \texttt{compare\_report\_id}; the imported publication profile, citation docket, carrier slot, audit trail, and escalation ladder remain auditable support for the normal form and reopen only when exact-location, full-audit, or reveal-order is itself the live question.

The maintained worked ledger now ships that shorter stop explicitly as \nolinkurl{answer_side_terminal_citation_normal_form} together with its one-clause \nolinkurl{answer_side_terminal_sentence_normal_form}. The same two fields are mirrored verbatim through \nolinkurl{example_question_routes.json} and the series spine, so a later terse note can quote the maintained answer-side stop directly instead of reconstructing it from the longer route tuple.

So the worked example now says not only what the answer-side companions are and in what order they may be revealed, but also which one a referee should receive for one concrete follow-up request. Each review entry now also exports an explicit \nolinkurl{review_request_class_vocabulary_mode} posture tag, and the worked route/spine surfaces mirror that through a shared \nolinkurl{answer_review_request_class_mode_map} so the repeated six request-class labels stay visibly menu-owned rather than being silently flattened into later route prose. Those same route/spine surfaces also mirror a new \nolinkurl{answer_escalation_step_mode_map} imported from the escalation ladder so the repeated five reveal-step labels stay visibly ladder-owned rather than being silently re-owned by the menu or its downstream mirrors. Its route catalog, \nolinkurl{example_question_routes.json}, now imports that six-class handoff map verbatim for each shipped audience/question route, together with the mirrored answer-side terminal-form fields, so a terse later note may cite one maintained route object for the exact brief/expanded/exact/full-audit/request-packet/reveal-order handoff choices while the review menu remains the canonical owner of those choices. The companion theorem-to-review bridge in Synthesis~55 then lets later notes connect this request-class menu back to the relevant line-item and mathematical roots without rewriting that lineage in prose.\cite{series:seriesspines}

\begin{thebibliography}{99}\small

\bibitem{series:theoremtopublicationspine}
Anonymous.
\newblock Theorem-to-Publication Spines for Terse Anonymity Papers: Earliest Sufficient Stops from Mathematics to Lineaged Release.
\newblock Series synthesis note (Synthesis~76), 2026. (This archive.)

\bibitem{series:challengeanswercitationdockets}
Anonymous.
\newblock Challenge Answer Citation Dockets for Successor Maintenance: Smallest-Sufficient Late-Stage Citation Choices from Capsules to Audit Trails.
\newblock Series synthesis note (Synthesis~50), 2026. (This archive.)

\bibitem{series:challengeanswerpublicationprofiles}
Anonymous.
\newblock Challenge Answer Publication Profiles for Successor Maintenance: Budgeted Reuse Modes from Answer Cards, Sentence Locks, and Citation Dockets.
\newblock Series synthesis note (Synthesis~51), 2026. (This archive.)

\bibitem{series:challengeanswerdisclosurepackets}
Anonymous.
\newblock Challenge Answer Disclosure Packets for Successor Maintenance: Public-on-Request Audit Slices from Publication Profiles, Audit Trails, and Support Manifests.
\newblock Series synthesis note (Synthesis~52), 2026. (This archive.)

\bibitem{series:challengeanswerescalationladders}
Anonymous.
\newblock Challenge Answer Escalation Ladders for Successor Maintenance: Ordered Brief-to-Audit Reveal Policies from Publication Profiles and Disclosure Packets.
\newblock Series synthesis note (Synthesis~53), 2026. (This archive.)

\bibitem{series:downstreamnormalforms}
Anonymous.
\newblock Downstream Normal Forms for Successor Maintenance: Typed Resolution and Exact Surface Assembly.
\newblock Series synthesis note (Synthesis~31), 2026. (This archive.)

\bibitem{series:releasespine}
Anonymous.
\newblock Release Spine Contracts for Successor Maintenance: Stable Stages and Stage-Local Handoffs.
\newblock Series synthesis note (Synthesis~32), 2026. (This archive.)

\bibitem{series:seriesspines}
Anonymous.
\newblock Series Spines from Mathematical Roots to Referee Handoffs: A Non-Redundant Bridge from Backbones, Receipts, Replay, Release Evolution, and Review Menus.
\newblock Series synthesis note (Synthesis~55), 2026. (This archive.)

\bibitem{series:workedexample}
Anonymous.
\newblock Worked Example: Receipt Interlock for an Anonymous-DHT Reader-Privacy Claim.
\newblock Series synthesis note (Synthesis~17), 2026. (This archive.)

\bibitem{series:publicrequestcontracts}
Anonymous.
\newblock Public Request Contracts for Source-Only Releases: Exact Paper-Level Handoffs from Boilerplate Policy to Maintained Packet Cuts.
\newblock Series synthesis note (Synthesis~56), 2026. (This archive.)

\bibitem{series:requestableevidenceclasses}
Anonymous.
\newblock Requestable Evidence Classes for Source-Only Releases: Stable Names for Omitted-Support Families behind Packet Cuts and Public Request Contracts.
\newblock Series synthesis note (Synthesis~57), 2026. (This archive.)

\bibitem{series:requestfulfillmentcertificates}
Anonymous.
\newblock Request Fulfillment Certificates for Source-Only Releases: Exact Completion Criteria for Public Request Contracts and Typed Omitted-Support Cuts.
\newblock Series synthesis note (Synthesis~58), 2026. (This archive.)



\bibitem{series:publicrequestresponsepacketrefreshresponsepacketclosureverdicts}
Anonymous.
\newblock Public Request Response Packet Refresh Response Packet Closure Verdicts for Source-Only Successor Releases: Stable Request-Side Terminal Forms after Visible Refresh.
\newblock Series synthesis note (Synthesis~73), 2026. (This archive.)

\bibitem{series:pairedterminalcitationdiscipline}
Anonymous.
\newblock Paired Terminal Citation Discipline for Review Spines: Answer-First Two-Anchor Reuse without Middle-Owner Drift.
\newblock Series synthesis note (Synthesis~74), 2026. (This archive.)

\bibitem{series:pairedterminalcutoverrules}
Anonymous.
\newblock Paired Terminal Cutover Rules for Review Spines: When Later Terse Papers Stop, Widen, or Reopen.
\newblock Series synthesis note (Synthesis~75), 2026. (This archive.)

\end{thebibliography}
\end{document}
